1.
===========================================================
Exercise 1.1
===========================================================
{true}                                 // GIVEN PRECONDITION
if a>=b && a>=c then                   // CONDITIONAL
    {true && a>=b && a>=c}
    -> 
    {a>=a && a>=b && a>=c}          // ASSIGNMENT
    max := a;
    {max>=a && max>=b && max>=c}       
else
    {true && !(a>=b && a>=c)} = {a<b || a<c}
    if b>=a && b>=c then               // CONDTIONAL
        {(a<b || a<c) && b>=a && b>=c}
        -> 
        {b>=a && b>=b && b>=c}      // ASSIGNMENT
        max := b;
        {max>=a && max>=b ∧ max>=c}
    else
        {(a<b || a<c) && !(b>=a && b>=c)}
        -> 
        {c>=a && c>=b && c>=c}      // ASSIGNMENT
        max := c;
        {max>=a && max>=b ∧ max>=c}
{max ≥ a ∧ max ≥ b ∧ max ≥ c}          // GIVEN POSCONDITION

===========================================================
Exercise 1.2
===========================================================
O triplo do exercício 1.1 é impreciso porque max pode ser qualquer valor desde que seja maior que a, b e c.
Solução: max tem de ser ou a, ou b ou c.

Triplo: {true} P {(max=a || max=b || max=c) && max>=a && max>=b && max>=c}

2.
===========================================================
Exercise 2.1
===========================================================
{n>=1}                                  // GIVEN PREDCONDITION
->                                  
9 = 10-1 && 1<=n                        
i := 1;
9 = (10^i)-1 && i<=n                    // ASSIGNMENT
r := 9;
r = (10^i)-1 && i<=n                    // INVARIANT
while i < n do
    {r = (10^i)-1 && i<=n && i<n}       // INVARIANT && GUARD
    -> 
    {r*10+9 = (10^(i+1))-1 && i+1<=n}   
    r := r*10+9;
    {r = (10^(i+1))-1 && i+1<=n}        // ASSIGNMENT
    i := i+1;
    {r = (10^i)-1 && i<=n}              // INVARIANT
{r = (10^i)-1 && i<=n && !(i<n)}        // INVARIANT && !GUARD
-> 
{r = (10^n)-1}

===========================================================
Exercise 2.2
===========================================================
O triplo continua válido. Pelo exercício anterior, provou-se que o triplo {n>=1} P {r=(10^n)-1} é válido. A condição mais forte (n>1) => 
a condição mais fraca (n>=1), já que > implica >=.
Assim, o conjunto de valores que satisfazem a condição mais forte (n>1) é um subconjunto do conjunto de valores que satisfazem a condição mais fraca (n>=1).
Pela left strengthening rule (φ → φ′ {φ′} S {ψ} LS | {φ} S {ψ}), com φ sendo a pré-condição mais forte e φ' a mais fraca, prova-se que o triplo continua 
válido com a pré-condição n>1.

===========================================================
Exercise 2.3
===========================================================
(1) n-i is a variant
{r = (10^i)-1 && i<=n && i<n && n-i = v} // INVARIANT && GUARD && v=e
->
{n-(i+1) < v}                            //ASSIGNMENT
r := r*10+9;
{n-(i+1) < v}                            //ASSIGNMENT
i := i+1;
{n-i < v}                                // v > e

(2) n-i não sai dos números naturais
r = (10^i)-1 && i<=n -> n-i>=0  				//INVARIANT -> e >= 0

3.
===========================================================
Exercise 3
===========================================================
O forward axiom tem um problema quando x ocorre em e: a pós-condição x=e nunca substitui e,
deixando-o para ser avaliado no pós-estado, exatamente como x. Isto não é correto, porque o
valor de e usado para calcular o novo x foi determinado no pré-estado, não no pós-estado, o que pode levar a incoenrências, como
no caso considerado. Neste caso, esta discrepância pode tornar a pós-condição contraditória, como
por exemplo, com x:=x+1, em que a pós-condição x=e torna-se x=x+1, avaliada inteiramente com o novo
valor de x — e nenhum valor de x satisfaz esta igualdade, logo o triplo {true} x:=x+1 {x=x+1}
não é válido, apesar de ser uma instância do forward axiom. O backward axiom conserta isto substituindo e por x, na pós-condição Q, formando a pré-condição Q[e/x], avaliada no 
pré-estado, tal como e foi de facto avaliado pela atribuição.
Isto garante que não há desencontro de estados, e portanto o axioma é correto mesmo quando
x ocorre em e.